\section{Wider guards and alternative comparison classes}
\label{sec:bb-guarded}

The transport lemmas also accept tables which differ on tapes never
visited by the initialized computation.  Those differences matter:
they change the complete source and image cylinders on which the
smooth extension is prescribed.  We give the remaining tables with
their actual guards and then isolate the extra pressure arguments.

\subsection{A shuttle aligned at the origin}
\label{sec:bb-shuttle-periodic}

\begin{theorem}\label{thm:bb-shuttle-periodic}
On the unit torus there is an effective smooth mean-zero force,
periodic from time zero with period one and with bounded mixed
derivatives, whose unique zero-data classical solution has
\[
 \exists t\ge0:\quad X_1(t;0)\pmod1\in(1/16,1/4)
 \quad\Longleftrightarrow\quad\text{halting}.
\]
The constructed pressure is zero.  Alternatively the force may be
projected to a solenoidal force with the same velocity and event,
and with the correspondingly changed mean-zero pressure.
\end{theorem}
\begin{proof}\label{proof:bb-shuttle-periodic}
Add a fresh nonhalting initial state $*$, without incoming transitions,
which writes back the scanned symbol and stays while entering the
original initial state.  This handles an originally halted input.
Use work, flag and history tracks
$\Sigma=\Gamma\times\{0,1\}\times(\{\perp,E\}\sqcup\mathcal R)$,
and controls $R_q,A_r,B_r,C_q,D_q$.  Its complete table is
\begin{equation}\label{eq:bb-shuttle-table}
\begin{array}{lll}
 R_q:(a,0,g)&\mapsto A_r:(b_r,0,g),d_r,\\
 A_r:(a,0,g)&\mapsto B_r:(a,1,g),1,\\
 B_r:(a,0,g)&\mapsto B_r:(a,0,g),1&g\ne E,\\
 B_r:(a,0,E)&\mapsto C_{p_r}:(a,0,r),1,\\
 C_q:(a,0,\perp)&\mapsto D_q:(a,0,E),-1,\\
 D_q:(a,0,g)&\mapsto D_q:(a,0,g),-1&g\ne E,\\
 D_q:(a,1,g)&\mapsto R_q:(a,0,g),0&g\ne E.
\end{array}
\end{equation}
In particular its first two rows do not exclude the frontier symbol;
this is a wider domain than the move-first table.  Into $A_r$ the
record identifies the work instruction; into $B_r$ the output flag
distinguishes entry and loop; into $C_q$ the record determines the
incoming scan; into $D_q$ a written $E$ distinguishes entry and loop;
into $R_q$ only clearing a flag is possible.  The destination fixes
the incoming displacement.  These checks include every copied track
and every permitted frontier case.

Initialize at $R_*$, with all flags zero and only $E$ at cell two.
At ready step $n$ the frontier is $e_n=n+2$, the records are at
$2,\ldots,n+1$, and the head is $h_n$, $|h_n|\le n$.
The work move reaches $h'=h_{n+1}<e_n$; the scan marks that cell,
records at $e_n$, creates the next frontier, returns to the mark and
clears it.  Its count is $2(e_n-h')+4$ and all its tests hold.
Thus it has exactly the required initialized behavior despite its
wider full domain.

Use even digits in base $b=2|\Sigma|$ and initial fractions
$\ell_*,r_*$.  Let $S$ be the control set and $m$ count the branches in
Lemma~\ref{lem:bb-onecell}, and put $L=100(1+|S|+m)$.
Give $R_*$ offset $-\ell_*$, other nonhalt states offsets $-4i$,
and ready halt states offsets $L+4i$.  The code is
$(o_s+\ell,r-r_*,0)$, so the initial point is zero.
The incoming conditions give separately disjoint full rectangle
families.  Thicken each by $[-1,1]$, lift to height $10i$, and use
the linear reciprocal middle scale.  Choose padding at most $1/10$
small enough for both family gaps.  Vertical phases can use fixed
column potentials $(0,\dot Z x\chi,0)$, with columns covering
$[-3,10i+3]$; middle phases use horizontal Hamiltonians in disjoint
height bands.  All paths and supports are contained in
$(-2L,2L)^3$: their coordinate bounds are
$[-4|S|-3,L+4|S|+4]$, $[-4,4]$, and $[-3,10m+3]$.
This verifies the hypotheses of the lift schedule on neighborhoods
of the whole boxes.

Scale by $\gamma=(8L)^{-1}$ into the torus patch
$(-1/2,1/2)^3$: $U(t,x)=\gamma W(t,x/\gamma)$.
The support lies inside $(-1/4,1/4)^3$, the potentials scale by
$\gamma^2$, and the residual is recomputed at viscosity $\nu$ by
Lemma~\ref{lem:p2-chart}.  Periodicity starts at zero
because the input is already encoded there and the pulse has idle
endpoint collars.  At integer times the particle is the correct code.
Halt first coordinates before scaling lie between $L+4$ and $2L$.
For a nonhalt branch its endpoint first coordinates are at most one;
the linear-scale identity keeps them at most one, and the weaker
bound $3/2$ also follows from center and width estimates.  After
scaling every nonhalt path lies in $(-1/4,1/16)$, disjoint modulo
one from the event.  The initialized dummy step makes the implication
valid even for an originally halted machine.

For the optional projection solve $\Delta q=\operatorname{div}g$,
$\int q=0$, where $g=\mathcal F_\nu[U]$.
Proposition~\ref{prop:projection-l2} gives the effective force
$f=g-\nabla q$ and pressure $p=-q$, with all mixed derivative bounds
and periodicity retained.  The pressure change is part of this
alternative; the direct residual had pressure zero.
\end{proof}

\subsection{A persistent marker and cursor shifts of two cells}
\label{sec:bb-marked-window}

\begin{theorem}\label{thm:bb-marked-window}
There is a unit-torus realization, periodic for $t\ge1$, with fixed
label $a_*=(1/4,1/4,1/4)$.  Its event
\[ \exists t\ge0:\quad X_1(t;a_*)\pmod1\in(2/3,5/6) \]
occurs exactly when the machine halts.
It uses an injective marked-head table whose cursor shifts may equal
two.  Its force is effective, smooth and mean zero, with bounded
mixed derivatives and the zero-data uniqueness class of
Proposition~\ref{prop:bb-realization}.
\end{theorem}
\begin{proof}\label{proof:bb-marked-window}
Normalize to one halt state and use work, permanent work-head marker,
and history tracks $(a,m,g)\in\mathcal A:=\Gamma\times\{0,1\}\times
(\{\bot,\#\}\sqcup\mathcal R)$.
Let $S$ be the set of controls $M_q,L_q,R_r$.  With the cursor at zero, make
exactly the following tests and writes, then reindex by
$T'_i=\bar T_{i+e}$ for the displayed shift $e$.
\begin{enumerate}
\item In $M_q$, $q$ nonhalting, require $m_0=1$ and, if $d_r\ne0$,
$m_{d_r}=0$.  Write $b_r$ at zero, move the marker to $d_r$ if
necessary, and enter $R_r$ with shift $e=d_r+1$.
\item In $R_r$, a cell with $m_0=0,g_0\ne\#$ gives shift $1$
in the same control.  If $m_0=0,g_0=\#,g_1=\bot$, write
$g_0=r,g_1=\#$ and enter $L_{p_r}$ with shift zero.
\item In $L_q$, if $m_0=0,g_0\ne\#$, shift $-1$ in that control.
If $m_0=1$, enter $M_q$ with shift zero.
\end{enumerate}
All unmentioned entries are retained.  Into $R_r$, a main entry has
output $m_{-1}=1$, whereas a scan has $m_{-1}=0$.
The record recovers the overwritten work letter, displacement and
marker changes.  Into $L_q$, a log write has output $g_1=\#$
and record $r$ at zero, whereas a scan has $g_1\ne\#$.
Its tests recover the overwritten pair.  Into $M_q$ only the
matching control change is possible.  This is a full-domain inverse;
uniqueness of a marker was not assumed on arbitrary tapes.

Initially the sole marker is at zero, the sole pointer at two, and
all other history is blank.  At work step $n$, head and marker are
at $h$, pointer at $n+2$, and records at $2,\ldots,n+1$.
The work update puts the marker at $h'=h+d_r$ and the cursor at
$h'+1\le n+2$.  The forward scan reaches the pointer, logs, and
returns to $h'$.  All tests hold since $h'<n+2$.  The cycle length is
\[
 1+[n+2-(h'+1)]+1+[(n+2)-h']+1.
\]
This proves finite continuation, exact work simulation and halting.

Use even digits, blank zero, in base $B=2|\mathcal A|+1$.
Enumerate every full word on $[-2,2]$ with an applicable rule, giving
$J$ branches.
Lemma~\ref{lem:bb-window-rectangles} gives prefix lengths $(2,3)$
and $(2+e,3-e)$, with $e\in\{-1,0,1,2\}$.
Each output length is at most four, so gaps after chart scaling
by $l=[100(|S|+1)]^{-1}$ are at least $lB^{-4}$.
Place the halt square at $(3/4,1/4)+l[0,1]^2$ and other squares
at $(1/4+2il,1/4)+l[0,1]^2$; the latter lie below $x_1=1/3$.
At base height $1/4$, lift branch $j$ to
$1/2+j/[4(J+1)]$, use exponential scales $B^{-e_j\theta},B^{e_j\theta}$,
and lower.  Padding
$\min(1/200,1/[32(J+1)],lB^{-4}/3)$ separates supports.
The middle Hamiltonian is
\[
 \dot c_x(y-c_y)-\dot c_y(x-c_x)
      -e_j\dot\theta\log B\,(x-c_x)(y-c_y).
\]
It gives the intended derivative on the full moving rectangle.
The plateau margins and finite scale bounds extend this to an open
neighborhood, as in Lemma~\ref{lem:bb-curl}.

Load the fixed label along its segment to the rational initial code
in time one, then repeat the unit field.  Integer samples give the
full local execution.  In a nonhalt period the center is below $1/3$
and the width at most $l$, so the whole path remains below $1/2$;
loading also avoids the event.  Halt codes lie in $[3/4,3/4+l]$.
Thus every real time is accounted for.  The pointer site $n+2$,
record count $n$, and relative pointer offset recover the absolute
head position at main checkpoints.  The residual construction
proves the remaining assertions.
\end{proof}

\subsection{A split direction shuttle in the whole space}
\label{sec:bb-direction-shuttle}

\begin{theorem}\label{thm:bb-direction-shuttle}
On $\R^3$ an effective compactly supported force with bounded mixed
derivatives and period one after $t=1$ realizes the halting event
$X_1(t;(1,0,0))<0$.  Uniqueness holds with bounded $u,\nabla u$,
$u\in C H^2\cap C^1L^2$ and pressure modulo constants in $C H^1$
on each finite interval.  An alternative force is square-integrable
in space-time, with the same label, event, support and derivative
bounds.
\end{theorem}
\begin{proof}\label{proof:bb-direction-shuttle}
Use exactly the direction table of Lemma~\ref{lem:bb-direction}
without its passive track, renaming
the controls by $M\mapsto S$, $G\mapsto R$, $A\mapsto I$,
and $B\mapsto J$, and the frontier by $F\mapsto\#$.  This is a bijection
of its entire partial table, including the nonfrontier tests.
Write $\Sigma$ for the resulting three-track alphabet.
Initialize frontier at one and zero flags.  At work step $j$ the
frontier is $j+1$ and the cycle length is
\[2(j-h_j)+5.\]

Use odd digits in base $B=2|\Sigma|+1$, offsets $3,6,\ldots$
for nonhalt controls and $-3,-6,\ldots$ for halt main controls.
In this realization split \emph{every} one-cell rule by every left
symbol $c$, including the stationary and right-moving rules.
Its source is $I(c)\times I(a)$, and the three target rectangles,
for moves $1,-1,0$, are respectively
\[
 \varphi_b\varphi_c([0,1])\times[0,1],\quad
 [0,1]\times\varphi_c\varphi_b([0,1]),\quad I(c)\times I(b),
 \qquad \varphi_a(t)=(d(a)+t)/B.
\]
The maps themselves are the restrictions of
\eqref{eq:bb-one-cell-branches}.  Source separation follows from
$(c,a)$, target separation from the fixed incoming move and $(c,b)$.
Their gaps are at least $B^{-2}$, and their linear parts are
$\operatorname{diag}(B^{-d},B^d)$.

Thicken by $[-1,1]$, lift branch $i$ to $4i$, use exponential
middle scale and center interpolation, then descend.  Padding smaller
than $(4B^2)^{-1}$ is valid, by source gaps, height gaps, and target
gaps in the three phases.  A first localized translation takes
$(1,0,0)$ to the rational initial code.  Exact full-box transport
then gives the local execution at every subsequent integer time.
Nonhalt centers have first coordinate at least three and horizontal
width at most one, so every intermediate point stays positive;
the nonhalt loading segment does too.  A halt sample has negative
first coordinate, including initial halting after loading.
Proposition~\ref{prop:bb-realization} gives the force, pressure zero,
and uniqueness.  The decaying alternative follows from
Proposition~\ref{prop:bb-late-logclock} below, applied to this same
loaded template.
\end{proof}

\subsection{Two-cell history and a gradient-only pressure condition}
\label{sec:bb-guarded-history}

\begin{theorem}\label{thm:bb-guarded-history}
An effective smooth force on $\R^3$, periodic from time zero,
with one compact spatial support and bounded mixed derivatives,
realizes the fixed event
\[
 \exists t\ge0:\quad X(t;(2,0,0))\in(-2,-1)\times(0,1)\times(-1,1)
 \quad\Longleftrightarrow\quad\text{halting}.
\]
Velocity is unique in the class
$u\in C([0,T];H^4)\cap C^1([0,T];H^2)$,
$\nabla p\in C([0,T];L^2)$ for every finite $T$, allowing a
distributional pressure gradient.  Alternatively the force belongs
to $L^2([0,\infty);H^j)$ for every integer $j\ge0$, with all the
same support, event and comparison-class conclusions.
\end{theorem}
\begin{proof}\label{proof:bb-guarded-history}
Add a fresh initial state which writes back and stays, and replace
all originally terminal state-symbol pairs by stay transitions into
one fresh terminal state $q_h$.  Normalize the other missing
instructions similarly.  Let $r=(q,a)\mapsto(p_r,b_r,d_r)$ denote
the resulting nonterminal instructions, and let $\mathcal A$ be the
work alphabet. The full cell alphabet is
$\Gamma=\mathcal A\times\{0,1\}\times(\{e,p\}\sqcup\mathcal R)$.
Use its work, flag and log tracks $W,F,L$,
and controls $N_q,A_r,R_r,B_q$.  The complete rules are:
\begin{enumerate}
\item At $N_q$, $q\ne q_h$, write $W_0=b_r$, enter $A_r$,
and shift $d_r$, with no test on the other tracks.
\item At $A_r$, require $F_0=0$, set it to one, enter $R_r$
and shift one.
\item At $R_r$, if $F_0=0,L_0\ne p$, shift one in that control.
If $F_0=0,L_0=p,L_1=e$, write $L_0=r,L_1=p$, enter $B_{p_r}$
and shift minus one.
\item At $B_q$, if $F_0=0,L_1\ne p$, shift minus one in that
control.  If $F_0=1$, clear it and enter $N_q$ without moving.
\end{enumerate}
The first output control records the overwritten work pair and shift.
Into $R_r$, output flag one at $-1$ distinguishes entry from its
loop.  Into $B_q$, a log insertion has pointer at output index two
and record at one, whereas a return loop has a nonpointer at two;
the record and tests recover both overwritten log entries.  Into
$N_q$ only unmarking is possible.  Thus this map is injective on
its whole domain, including the unrestricted first-row tracks.

Initialize zero flags and a sole pointer at cell three, all other
log cells empty.  At work step $n$, records occupy $3,\ldots,n+2$
and pointer $g=n+3$.  The work move reaches $h'$ with $g-h'\ge2$.
The finite scan flags $h'$, logs at $g$, advances the pointer, and
returns to the flag.  The return guard holds because the new pointer
is at $g+1$.  The cycle has $2(g-h'-1)+4$ steps.  This verifies
the work invariant and eventual terminal entry.

Use odd digits in radix $D=2|\Gamma|+1$ and write
$\phi(\xi)=\sum_{j\ge1}d(\xi_j)D^{-j}$ for the radix sum.
Let $L_*,R_*$ be the initialized left and right streams, ordered as in
Lemma~\ref{lem:bb-onecell}.  Give the terminal state
chart offset $(-2,0)$, the initial chart offset
$(2-\phi(L_*),-\phi(R_*))$, and all other charts $(4+2j,0)$.
The initial code is $(2,0,0)$, every nonterminal interval lies in
$x_1>1$, and terminal codes lie strictly in $(-2,-1)\times(0,1)$.
The infinite tails are in $[\delta,1-\delta]$,
$\delta=1/(D-1)$, so finite-prefix decoding of length $b$ is possible
from error less than $(\delta/2)D^{-b}$ without a boundary ambiguity.

Split full triples $[-1,1]$.  By
Lemma~\ref{lem:bb-window-rectangles}, source prefix lengths $(1,2)$
become $(1+d,2-d)$, target cylinders are complete, and both rectangle
families have gaps at least $D^{-3}$.  Thicken by $[-1/4,1/4]$,
lift to height $4i$, use linear middle scale
$\mu_i=1+(D^{-d_i}-1)\theta$, and descend.  Padding smaller than
$D^{-3}/4$ separates all phases.  Repeating the pulse, zero on
$[0,1/12]$ and $[11/12,1]$, gives periodicity from zero without
loading.  At a nonhalting branch's lift and descent the first
coordinate is greater than one; during the middle transfer the
particle's third coordinate is exactly $4i$, outside $(-1,1)$.
Thus even a horizontal route passing over the observation box cannot
cause a false visit.  The integer samples and the compiler invariant
prove halting equivalence.  The record count and pointer site $n+3$
recover absolute indexing at checkpoints.

For uniqueness apply Lemma~\ref{lem:b-comparison}(iii).
The explicit compactly supported smooth field lies in $CH^4\cap C^1H^2$
on every finite interval and has pressure zero.  The competing class
in the theorem is precisely case (iii), including its distributional
$CL^2$ gradient.  Lemma~\ref{lem:b-gradient} cancels that gradient
against the solenoidal velocity difference without introducing a pressure
$L^2$ norm.  Thus the velocity and pressure gradient are unique, while
pressure itself remains free up to a time-dependent spatial constant.
Proposition~\ref{prop:bb-late-logclock} gives the stated
decaying alternative.
\end{proof}

\subsection{A flagged head and two-pass parking}
\label{sec:bb-two-pass-parking}

\begin{theorem}\label{thm:bb-two-pass-parking}
For any machine and finitely specified initial tape, an effective
smooth force on $\R^3$, periodic from time zero, with one compact
spatial support and bounded mixed derivatives, realizes the event
\[
 \exists t\ge0:\quad X(t;0)\in(-6,-4)\times(-1,1)\times(-1,1)
 \quad\Longleftrightarrow\quad\text{halting}.
\]
The zero-data velocity is unique among classical smooth solutions
with bounded $u,\nabla u$, $u\in C H^2\cap C^1L^2$ and $p\in C L^2$
on every finite interval.  The prescribed solution has uniformly
bounded kinetic energy.
\end{theorem}
\begin{proof}\label{proof:bb-two-pass-parking}
Merge terminal states into $q_\dagger$, send missing instructions
there by a stay step, and add a fresh initial state $q_*$, with no
incoming instruction, which writes back and stays while entering
the old initial state.  For each normalized nonterminal rule
$r=(q,a)\mapsto(p_r,b_r,d_r)$ use cell letters
\[
 (w,h,j)\in\Gamma:=\mathcal A\times
 (\{\circ,\star\}\sqcup\{[r]:r\in\mathcal R\})\times\{0,1\},
\]
where $\mathcal A$ is the work alphabet and $\mathcal R$ is the
finite set of normalized instruction labels. Thus $\Gamma$ includes
all three tracks.
The complete table, with $h\ne\star$ wherever $h$ occurs, is
\[
\begin{array}{lll}
 M_q:(a,h,1)&\mapsto A_r:(b_r,h,0),d_r,\\
 A_r:(w,h,0)&\mapsto R_r:(w,h,1),1,\\
 R_r:(w,h,0)&\mapsto R_r:(w,h,0),1,\\
 R_r:(w,\star,0)&\mapsto B_{p_r}:(w,[r],0),1,\\
 B_q:(w,\circ,0)&\mapsto L_q:(w,\star,0),-1,\\
 L_q:(w,h,0)&\mapsto L_q:(w,h,0),-1,\\
 L_q:(w,h,1)&\mapsto M_q:(w,h,1),0.
\end{array}
\]
The destination fixes the incoming displacement.  The record index
recovers the first row's old work pair; a flag distinguishes the
entries into $R_r$; a record distinguishes the entries into $B_q$;
a written frontier distinguishes entry and loop into $L_q$;
and entry into $M_q$ is the unique cell-preserving last row.
This proves the full one-cell inverse conditions.

Initialize the sole flag at zero and the sole frontier at three,
with empty history elsewhere.  At ready step $k$, only the work
head $i$ is flagged, records occupy $3,\ldots,k+2$, and the frontier
is $e_k=k+3$.  The work rule clears the flag and moves to
$j=i+d_r\le e_k-2$, which is unflagged even for $d_r=0$.
It is flagged anew; the scan logs at $e_k$, advances the pointer,
returns, and preserves the ready flag.  The two loop lengths are
$e_k-j-1$ and $e_k-j$, with five other rules, for a total
$2(e_k-j)+4$.  Thus all scans terminate and ready halting is exact.

Use odd digits in base $K=2|\Gamma|+2$, initial fractions $x_*,y_*$,
and state offsets $a_{M_{q_*}}=0$, $a_{M_{q_\dagger}}=-5$, and
$4,8,12,\ldots$ for the others.  The code
$(a_s-x_*+x,y-y_*,0)$ starts at zero.  Lemma~\ref{lem:bb-onecell}
gives $N$ full rectangle maps with scales $\lambda_i\in\{K^{-1},1,K\}$;
thicken by $[-1/4,1/4]$.  The entire terminal block is inside the
observation box because $0<x_*,y_*<1$.  Every nonterminal block is
in $x_1>-1$.  Eventually constant tails are rational even for an
arbitrary finitely specified initial tape, and the pointer and record
count recover its absolute head position at ready checkpoints.

For the required two-pass geometry, let $g$ be the minimum of one
and all positive separating coordinate gaps in either planar family,
and set $\epsilon=g/20$.  Choose
\[
 C=20+\max(\{0\}\cup\{\text{right endpoints of source and target boxes}\}),
 \qquad S_i=(C+4i,0,0).
\]
In a first pass give each source four consecutive slots: lift by ten,
translate above $S_i$, descend, and reshape there by the linear
reciprocal scale.  In a second pass give each parked box three slots:
lift by ten, translate above its target, and descend.  The total is
$7N$ slots, and every source is gone before a target is filled.
Widths remain at most one.  Source gaps protect the first lift,
parking spacing protects strain and parking moves, target gaps protect
delivery, and travel height separates every horizontal transfer from
resting boxes.  These inequalities remain strict under $2\epsilon$
padding.  Applying Lemma~\ref{lem:bb-curl} in chronological order
fixes every stationary box and gives the exact full-box map to the
moving one.  Thus overlapping source and target families cause no
interference.

Repeat the flat-ended template from zero and take its residual.
The integer sampling induction gives the table execution.  In a
nonhalting run every endpoint box lies in $x_1>-1$; translations
interpolate corresponding endpoint points, strain occurs far to the
right, and waiting leaves the point fixed.  The entire two-pass path
therefore remains in $x_1>-1$.  A terminal sample lies inside the
observer.  This proves the event for all times.

For the exact pressure class use Lemma~\ref{lem:b-comparison}(iv).
The prescribed compactly supported smooth velocity satisfies the
reference $CH^2\cap C^1L^2$ and bounded-gradient conditions.  The
competing velocity has precisely these bounds and a $CL^2$ pressure
representative, as that case requires.  Its cutoff argument uses the
pressure itself, with error $CR^{-1}\|p\|_2\|u-U\|_2$, and assumes
no extra integrability of its derivative.  The lemma gives velocity
uniqueness; the $L^2$ pressure representative is then zero.
Compact support and
bounded derivatives of the prescribed velocity give its uniform
energy bound and its global material flow.
\end{proof}


\subsection{A two-cell append and its guarded inverse}\label{bc:two-cell-section}
Use Lemma~\ref{lem:bb-window-compiler} with initial frontier $r_0=1$,
renaming Ready, Go and Back as $M,R,L$ and $(E,P)$ as
$(\perp,\star)$. The track order is work, mark, history.
The Ready history coordinate and the marked Back row remain unrestricted;
the unmarked Back row tests precisely the right-neighbor history.
Thus the full table, not only its initialized execution, is identical.
At checkpoint $n$, its frontier is $p_n=1+n>h_n$, where $h_n$ is
the work-head position, and the exact cycle length is $2(1+n-h_n)+1$.
Lemma~\ref{lem:bb-window-rectangles} applies with
$(a,b)=(1,2)$: the complete allowed triple fixes both history writes,
and its full image has prefix lengths $(1+e,2-e)$.

\begin{theorem}[Guarded cylinders in a torus chart]\label{bc:guarded-torus}
On the unit torus, the fixed label $(1/8,0,0)$ and the strip
$\{x:x_1\pmod1\in(-1/4,0)\}$ realize halting by a general mean-zero
force with zero initial velocity and pressure. It is periodic after time
one and all its mixed derivatives are bounded. The torus comparison
class gives uniqueness.
\end{theorem}
\begin{proof}\label{bc:guarded-torus-proof}
Let $a$ be the augmented alphabet size and $k$ the control count in
Lemma~\ref{lem:bb-window-compiler}. Use zero-blank base $B=2a$, and enumerate
nonterminal and terminal states with positive integer indices $j\ge1$
and first-coordinate offsets respectively
$L_0+3j$ and $-L_0-3j$, where $L_0=4k+ka^3+10$.
Refinement by triples gives at most $N=ka^3$ branches. Their source
prefix lengths are $(1,2)$ and image lengths $(1+e,2-e)$, so the
smallest possible grid gap is $B^{-3}$. If $P_i,Q_i$ are lower corners,
the branch map is
$Q_i+\diag(B^{-e_i},B^{e_i})(y-P_i)$.

Choose three flat ramps $\alpha,\beta,\gamma$, active respectively in
$(1/8,1/4)$, $(3/8,5/8)$ and $(3/4,7/8)$. Set
\[
 c_i(t)=((1-\beta)P_i+\beta Q_i,\ i(\alpha-\gamma)),
 \quad l_i=1-\beta+\beta B^{-e_i}.
\]
The moving rectangle has side lengths $l_i B^{-1}$ and
$l_i^{-1}B^{-2}$. Use a physical padding
$\varepsilon=B^{-3}/4$ around each moving rectangle in three dimensions.
During ascent and descent the planar endpoint gaps separate these
neighborhoods; during the affine phase the unit height gaps separate them.
Use Lemma~\ref{lem:bb-curl} for the paths
$c_i+\diag(l_i,l_i^{-1},1)(y-(P_i,0))$ and sum the disjoint fields.
The first coordinate is an endpoint convex combination even though the
paths are based at lower corners rather than centers. All supports lie
in $[-2L_0,2L_0]^3$, by the definitions of $L_0$ and $N$.

Scale by $\rho=1/(32L_0)$, so the repeated mechanism is supported in
$[-1/16,1/16]^3$. Load the fixed label $(1/8,0,0)$ to the scaled
rational input code by a moving translation with cutoff radius $1/32$;
its support stays in $(-1/4,1/4)^3$. For a nonterminal initial code,
the entire first coordinate during loading is positive, and subsequent
nonterminal branches have first coordinate greater than $1/32$.
Terminal codes have first coordinate less than $-1/32$. Hence the stated
strip detects exactly terminal entry, including initial halting, without
any wraparound artifact. Periodize the compact curl fields and recompute
the force at physical $\nu$. Mean zero, bounded derivatives and comparison
follow from Lemma~\ref{bc:analytic}.
\end{proof}

\subsection{A filled half-line and mean-corrected vertical motion}
\label{bc:halfline-section}
The next recorder starts with an infinite constant history tail rather
than a single frontier symbol. Its guard and its mean-zero vertical
realization both have independent roles.

\begin{lemma}[History at a filled half-line]\label{bc:halfline-table}
A finite guarded table has a full-domain inverse and simulates the original
machine from a tape with a separately constant tail in each direction.
At checkpoint $n$ its nonempty history cells are exactly the absolute
indices below $2+n$.
\end{lemma}
\begin{proof}\label{bc:halfline-table-proof}
Normalize to one halt state $h$. For $i=(q,a)\in I=(Q\setminus\{h\})\times A$,
write $\delta(q,a)=(q_i,b_i,e_i)$. History letters are
$\mathcal H=\{\perp,\#\}\sqcup I$, with
$\mathcal H_+=\{\#\}\sqcup I$. The tracks are work, mark and history;
controls are $\mathsf R_q,\mathsf B_q,\mathsf T_i,\mathsf S_i$.
Write $\ell_j$ for the history at relative cell $j$.
Here is the complete table, with unmentioned cells unchanged:
\begin{equation}\label{bc:halfline-rules}
\begin{array}{c|c|c|l}
\text{input}&\text{output}&\text{move}&\text{condition}\\\hline
\mathsf R_q(a,0,k)&\mathsf T_i(b_i,0,k)&e_i&i=(q,a),\ q\ne h\\
\mathsf T_i(c,0,k)&\mathsf S_i(c,1,k)&1&k\in\mathcal H_+\\
\mathsf S_i(c,0,k)&\mathsf S_i(c,0,k)&1&k\in\mathcal H_+\\
\mathsf S_i(c,0,\perp)&\mathsf B_{q_i}(c,0,i)&0&\ell_{+1}=\perp\\
\mathsf B_q(c,0,k)&\mathsf B_q(c,0,k)&-1&k\in\mathcal H_+\\
\mathsf B_q(c,1,k)&\mathsf R_q(c,0,k)&0&k\in\mathcal H_+.
\end{array}
\end{equation}
Initialize $\mathsf R_{q_0}$ at head zero with marks zero, history $\#$
at every index below two and $\perp$ elsewhere. At checkpoint $n$, the
work head $p_n$ has $|p_n|\le n$ and the boundary is $r_n=2+n$.
After the work move, $p'=p_n+e_i\le n+1<r_n$. Thus the mark is placed
on nonempty history, and the right scan reaches the first empty cell.
Its right neighbor is also empty, so the append is allowed. A finite
left scan returns to and erases the unique mark. The new boundary is
$r_n+1$; no other work symbol changes. The filler boundary at absolute
index two is permanent, so it also retains the tape's absolute alignment.

Into $\mathsf T_i$ the tag gives the old work instruction and shift.
Into $\mathsf S_i$, both shifts are rightward, and the mark at the old
cell distinguishes entry from a loop. Into $\mathsf B_q$, an output
with empty history at offset $+1$ must be an append: its current record
identifies $i$. An output of a leftward loop has nonempty history at
that offset, since it is the old current cell. This is exactly why the
append guard is necessary. Into $\mathsf R_q$ only mark erasure applies.
Each case gives a unique predecessor on the full domain.  In particular,
entry into $\mathsf B_q$ may have displacement zero or minus one;
the neighboring history, rather than the target control alone, recovers
that displacement.  Apply Lemma~\ref{lem:bb-window-rectangles} with
$(a,b)=(1,2)$, refining by the full cells at $-1,0,1$.
Its full-domain injectivity hypothesis has just been proved, and every
write and displacement lies in the allowed window.  It therefore gives
full image cylinders and separately disjoint closed rectangle families.
\end{proof}

Enumerate the resulting branches by $i=1,\ldots,N$, with displacements
$e_i$.  Number the augmented alphabet by $0,\ldots,k-1$, use base $b=3k+1$
and digits $1+3\ell$. Two different prefixes have gaps at least
$2b^{-r}$ at their first differing place $r$. Codes lie in the interior
of their intervals. If there are $K$ controls, let $\sigma_0=1/[8(K+1)]$.
Put the terminal state's lower corner at $(3/4,1/4)$ and the other lower
corners at $(1/8+2\ell\sigma_0,1/4)$ for $0\le\ell<K-1$, giving every state square side
$\sigma_0$. Their branch rectangles have sides
$\sigma_0b^{-1},\sigma_0b^{-2}$ and image sides
$\sigma_0b^{-(1+e)},\sigma_0b^{-(2-e)}$. The families are separately
disjoint by the full inverse. Every endpoint rectangle is contained in
\begin{equation}\label{bc:halfline-code}
 F=[1/8,7/8]\times[1/4,3/8],\qquad z_*=1/4.
\end{equation}
All nonterminal state squares have first coordinate at most $3/8$.

\begin{proposition}[Mean-corrected vertical transport]\label{bc:halfline-transport}
The preceding full rectangles at height $z_*$ have an effective mean-zero
solenoidal realization on the unit torus. Every tracked first coordinate
is between its endpoint values throughout the transition.
\end{proposition}
\begin{proof}\label{bc:halfline-transport-proof}
For a source rectangle $D$, choose a planar cutoff $\theta_D$ with
padding $\varepsilon=\sigma_0b^{-3}/4$, equal to one near $D$. Define
\[
 \kappa_D(x,y)=\theta_D(x,y)-\theta_D(x,y-1/2)
\]
periodically. It has integral zero. On its own source it equals one,
and on all other sources it equals zero; the shifted copy lies away from
$F$. Use the same construction for target rectangles $E$. Assign branch
$i$ the height $z_i=1/2+i/[4(N+1)]$ and choose disjoint height cutoffs
$\rho_i$ equal to one near $z_i$, with padding $1/[16(N+1)]$. The
vertical fields with speeds
\[
 Z_D(x,y)=\sum_i(z_i-z_*)\kappa_{D_i}(x,y),\qquad
 Z_E(x,y)=\sum_i(z_i-z_*)\kappa_{E_i}(x,y)
\]
are divergence free and mean zero, because they are independent of $z$.

Let $\chi$ equal one near $F$, with planar padding $1/32$. For branch
$i$, let $p_i,q_i$ be the source and target centers and put $d_i=q_i-p_i$, $c_i(\mu)=p_i+\mu d_i$,
$l_i=1-\mu+\mu b^{-e_i}$ and
\[
 H_i=\chi\left[d_{i,1}y-d_{i,2}x+
 \frac{b^{-e_i}-1}{l_i}(x-c_{i,1})(y-c_{i,2})\right].
\]
The middle field $W(\mu)=\sum_i\rho_i(z)(\partial_yH_i,-\partial_xH_i,0)$
is divergence free and mean zero. The path
\begin{equation}\label{bc:halfline-path}
 c_i(\mu)+\diag(l_i,l_i^{-1})(\zeta-p_i)
\end{equation}
stays in $F$: the first side is interpolated linearly and the reciprocal
inequality \eqref{eq:ba-reciprocal-envelope} bounds the second side
by the interpolated endpoint sides. Thus $\chi=1$ along the entire
path, and differentiation gives its exact velocity.
For three ordered flat ramps $\alpha_0,\alpha_1,\alpha_2$ on a unit
interval, the complete field is
\begin{equation}\label{bc:halfline-period}
 V=\alpha_0'(0,0,Z_D)+\alpha_1'W(\alpha_1)
                   -\alpha_2'(0,0,Z_E).
\end{equation}
It lifts each source, performs its affine map in its own height layer,
and lowers it at its target. The first coordinate is unchanged during
vertical motion and is an endpoint convex combination in the middle.
All statements hold on the full rectangles by the cutoff plateaus and
ODE uniqueness. The zero-mean correction acts away from the computational
rectangles and therefore changes none of these paths.
\end{proof}

\begin{theorem}[A fixed strip with half-line history]\label{bc:halfline-event}
On $\T^3$, the fixed label $(1/8,1/4,1/4)$ and strip
$2/3<x_1<15/16$ realize halting by a general mean-zero force, with zero
initial velocity and pressure, bounded mixed derivatives and period one
after time one. For a fixed machine only the loading field depends on the
input. The material maps are volume-preserving diffeomorphisms, and the
velocity and material-acceleration identities of
Lemma~\ref{lem:p2-realization} hold.
\end{theorem}
\begin{proof}\label{bc:halfline-event-proof}
The two initial tape tails are separately constant, so their radix sums
are computable rationals, even though the left history tail is filled.
Load the fixed label to that code at height $z_*$ by the planar Hamiltonian
$\chi(d_x y-d_y x)$ with a flat scalar schedule. The segment stays in $F$.
A nonterminal load and every nonterminal phase have first coordinate at
most $3/8$; a terminal code lies in the stated strip. For an initially
terminal machine the loading endpoint suffices. Repeat
\eqref{bc:halfline-period} after time one. Checkpoint simulation and the
intermediate-time guard prove the event equivalence. The common analytic
lemma proves force effectivity, uniqueness, boundedness and the asserted
flow identities. In particular $U=(\partial_t\Phi_t)\circ\Phi_t^{-1}$ and
$\partial_{tt}\Phi_t=(f-\nabla p+\nu\Delta U)\circ\Phi_t$; these are
identities for the constructed flow, not additional unknown equations.
\end{proof}
